\documentclass[11pt,a4paper]{amsart}
\usepackage[T1]{fontenc}
\usepackage{lmodern,microtype}
\usepackage[margin=24mm,headheight=14pt]{geometry}
\usepackage{amsmath,amssymb,amsthm,mathtools}
\usepackage[hidelinks]{hyperref}
\numberwithin{equation}{section}
\newtheorem{theorem}{Theorem}
\newtheorem{lemma}[theorem]{Lemma}
\newtheorem{proposition}[theorem]{Proposition}
\theoremstyle{remark}\newtheorem{remark}[theorem]{Remark}
\newcommand{\G}{\mathbb G}
\newcommand{\R}{\mathbb R}
\newcommand{\eps}{\varepsilon}
\newcommand{\tr}{\operatorname{tr}}
\newcommand{\Sym}{\operatorname{Sym}}
\newcommand{\ad}{\operatorname{ad}}
\newcommand{\cH}{\mathcal H}
\newcommand{\cA}{\mathcal A}
\newcommand{\cD}{\mathcal D}
\newcommand{\cV}{\mathcal V}
\newcommand{\cC}{\mathcal C}
\newcommand{\cU}{\mathcal U}
\newcommand{\cE}{\mathcal E}
\newcommand{\Mom}{\mathsf M}
\newcommand{\ip}[2]{\langle #1,#2\rangle}
\newcommand{\norm}[1]{\lVert #1\rVert}
\newcommand{\dd}{\,d}
\allowdisplaybreaks[1]
\title{A scale--tensor proof of an ABP estimate on stratified Carnot groups}
\date{October 6, 2026}
\begin{document}
\begin{abstract}
We develop a large-$p$ ABP estimate for arbitrary bounded measurable symmetric
horizontal coefficient matrices on a stratified Carnot group. The regularizer is
an all-layer, time-dependent left-invariant diffusion. Its exact tensor lift is
renormalized by its finite-dimensional zero-order moment evolution. The
normalized nonhorizontal blocks have zero scale-boundary data and positive
scale damping. Their energies at layer $k$ have constants depending only on
smoothing weights strictly below $k$. Two integrations by parts pair the
remaining fluxes with scalar gradient dissipation and yield an inverse current-
layer smoothing factor. A finite graded parameter selection, followed by a
small choice of $q-1$, closes the forward adjoint entropy inequality. The
elliptic maximum estimate follows by resolvent duality, a horizontal quadratic
barrier, and smooth approximation of the coefficients.
\end{abstract}
\maketitle

\section{The result and the role of the new correction}
Let $\mathfrak g=V_1\oplus\cdots\oplus V_s$ be a stratified Lie algebra,
$r=\dim V_1$, and $D=\sum_{j=1}^s j\dim V_j$. Fix orthonormal bases
$Z_{j,\alpha}$ of the layers and write $X_\alpha=Z_{1,\alpha}$. All fields
are left invariant. Haar measure is understood. Let
\[
 L_Au=\sum_{\alpha,\beta=1}^r a_{\alpha\beta}X_\alpha X_\beta u.
\]
Symmetry of $A$ means that this is contraction with the symmetrized horizontal
Hessian. A fixed left-invariant homogeneous distance defines the balls below.

\begin{theorem}\label{thm:abp}
For every stratified Carnot group $\G$ and $E\ge1$, there is a finite
$p_0=p_0(\G,E)>1$ with the following property. Let $\Omega\subset B_R(g_0)$
be bounded and open. Suppose $A:\Omega\to\Sym_r$ is Borel measurable,
$\lambda I\le A\le\Lambda I$ almost everywhere, and $E=\Lambda/\lambda$.
If $u\in C^2(\Omega)\cap C(\overline\Omega)$ and $L_Au\ge f$ almost
everywhere, then, for $p_0\le p\le\infty$,
\begin{equation}\label{eq:abp}
 \sup_\Omega u\le b+
 \frac{C(\G,E,p)}\lambda R^{2-D/p}
 \norm{f^-\mathbf1_{\{u>b\}}}_{L^p(\Omega)},
 \qquad b=\sup_{\partial\Omega}u^+.
\end{equation}
The constants do not involve coefficient derivatives or a modulus of
continuity. No lattice or boundary regularity assumption is required.
\end{theorem}

The common-flow entropy calculation comes from the supplied kinetic and
Heisenberg notes \cite{Kinetic,Heisenberg}. The additional mechanism here is to
remove the full tensor's zero-order moment evolution before estimating its
higher blocks. This is different from either discarding the tensor commutator
or requiring the complete regularized tensor to be pointwise elliptic.

Normalize $\lambda=R=1$ until the final section. Initially $A$ is smooth,
$I\le A\le EI$, and is constant outside a compact set. Let
\begin{equation}\label{eq:physical}
 m_t=L_A^*m=\sum_{\alpha,\beta}X_\alpha X_\beta(a_{\alpha\beta}m),
 \qquad m>0,\qquad\int m=1.
\end{equation}
The a priori estimates are uniform over these smooth coefficients. Smooth
positive data with sufficiently rapid polynomial decay can be used initially;
Section~\ref{sec:analytic} explains the analytic reduction and approximation.

\section{The scalar regularizer and the exact entropy rule}
Choose $a_1=1$ and $a_j\ge1$, $2\le j\le s$, to be fixed before choosing $q$.
Put
\begin{equation}\label{eq:scalegen}
 \cA_h=\frac12\sum_{j=1}^s a_jh^{j-1}\sum_\alpha Z_{j,\alpha}^2,
 \qquad \partial_h S_h=\cA_hS_h,\qquad S_0=I.
\end{equation}
The evolution is nonautonomous; $S_h$ is not being identified with
$\exp(h\cA_h)$. It is positive and preserves mass. It is a convolution
regularizer with a smooth positive kernel $K_h$, with convolution order fixed
by \eqref{eq:scalegen}. Homogeneous scaling gives
\begin{equation}\label{eq:kernel}
 K_h(g)=h^{-D/2}K_1(\delta_{h^{-1/2}}g),\qquad \int K_h=1,
 \qquad\norm{S_h}_{1\to\infty}\le C_\G h^{-D/2}.
\end{equation}
The last constant can be taken independent of the $a_j\ge1$. Indeed, the
horizontal Nash inequality and
\[
 \frac d{dh}\norm{S(h,h_0)v}_2^2
 =-\sum_j a_jh^{j-1}\norm{\nabla_jS(h,h_0)v}_2^2
 \le-\norm{\nabla_HS(h,h_0)v}_2^2
\]
give the $1\to2$ bound. Applying the same argument to the reverse-time adjoint
evolution gives $2\to\infty$; composition gives \eqref{eq:kernel}. The
horizontal Nash inequality follows from the fixed-sub-Laplacian heat bound by
$\norm{v}_2^2\le h\norm{\nabla_Hv}_2^2+C h^{-D/2}\norm{v}_1^2$ and
optimization. Only fixed-operator heat facts are used here; see \cite{GTHeat}.

At a fixed physical time, put
\[
 \rho=S_hm,\quad g=\log\rho,\quad q=1+\eps,\quad 0<\eps\le\tfrac12,
 \quad Q(t,h)=\int\rho^q.
\]
For a scalar or vector $v$, define
\begin{align}
 \cD_q(v)=\sum_{j,\alpha}a_jh^{j-1}\int\rho^q
 \bigg[&\frac{2-q}{q}|Z_{j,\alpha}v|^2\\[-2mm]
 &+\frac{q\eps}{2}\left|vZ_{j,\alpha}g+\frac2q Z_{j,\alpha}v\right|^2
 \bigg].\label{eq:Dq}
\end{align}
If $U_h=\cA_hU+F$ and $v=U/\rho$, then
\begin{equation}\label{eq:entropy-rule}
 \frac d{dh}\int\rho^{q-2}|U|^2
 =-\cD_q(v)+2\int\rho^{q-2}\ip UF.
\end{equation}
In addition to $\cD_q(v)\ge\frac13\sum a_jh^{j-1}\int\rho^q|Z_jv|^2$,
we will repeatedly use the exact second form
\begin{equation}\label{eq:second-square}
 \cD_q(v)=\sum_{j,\alpha}a_jh^{j-1}\int\rho^q
 \left[|Z_{j,\alpha}v+\eps vZ_{j,\alpha}g|^2
 +\frac{\eps(1-\eps)}2|v|^2(Z_{j,\alpha}g)^2\right].
\end{equation}
This is why no loss $1/(q-1)$ is needed in the source estimates below.

Write
\begin{align}
 I_j(h)&=\sum_\alpha\int\rho^q(Z_{j,\alpha}g)^2,&
 J_j&=\int_0^1h^{j-1}I_j(h)\dd h,\label{eq:J}\\
 L_j&=\sum_\alpha\int_0^1h^j\cD_q(Z_{j,\alpha}g)\dd h,&
 J_*&=\sum_jJ_j,\quad L=\sum_jL_j.\label{eq:L}
\end{align}
Then
\begin{equation}\label{eq:Qgap}
 Q(t,0)-Q(t,1)=\frac{q\eps}{2}\sum_ja_jJ_j.
\end{equation}

\section{Uniform scalar estimates and finite parameter selection}
All constants in this section are uniform for $1<q\le3/2$.

\begin{lemma}\label{lem:scalar}
Given finitely many polynomials $P_k(a_1,\ldots,a_{k-1})$ with nonnegative
coefficients, and $\delta>0$, the weights can be chosen so that
\begin{equation}\label{eq:scalar-targets}
 J_*\le2J_1,\qquad L\le C_\G J_1,\qquad
 \sum_{k=2}^sP_k(a_{<k})J_k\le\delta J_1.
\end{equation}
One may impose simultaneously any finite collection of conditions
$P_k(a_{<k})/a_k\le\delta$. The constant $C_\G$ does not depend on the
values of the weights selected in this way or on $q$.
\end{lemma}

\begin{proof}
For $U\in V_i$, $V\in V_j$, write $g_U=Ug$. In the entropy of $U\rho$, the
commutator with $a_jh^{j-1}V^2/2$ has source
\begin{align}
 -2a_jh^{j-1}\int\rho^q
   (Vg_U+\eps g_Vg_U)[U,V]g
 +a_jh^{j-1}\int\rho^q g_U[[U,V],V]g.\label{eq:scalar-source}
\end{align}
This follows from $[U,V^2]=2V[U,V]+[[U,V],V]$ and one integration by parts.
Use \eqref{eq:second-square}, multiply by $h^i$, integrate, and discard the
nonnegative upper endpoint. Summing the basis fields gives
\begin{equation}\label{eq:scalar-upper}
 L_i\le C_\G\left[iJ_i+
 \sum_{j:i+j\le s}a_jJ_{i+j}+
 \sum_{j:i+2j\le s}a_j\sqrt{J_iJ_{i+2j}}\right].
\end{equation}

Represent a basis of $V_k$ by fixed linear combinations of
$[X,Y]$, $X\in V_1$, $Y\in V_{k-1}$. The identity
\[
 \int\rho^q(Z_kg)^2
 =\sum c\int\rho^q\big[(Xg)Y(Z_kg)-(Yg)X(Z_kg)\big]
\]
has no derivative-of-weight remainder. Integrating with weight $h^{k-1}$ and
commuting $Y$ and $X$ to the right of $Z_k$ yields
\begin{equation}\label{eq:scalar-transfer}
 J_k\le C_\G\left[
 \sqrt{J_1\left(L_{k-1}/a_k+J_{2k-1}\right)}+
 \sqrt{J_{k-1}\left(L_1/a_k+J_{k+1}\right)}\right],
\end{equation}
where $J_l=0$ for $l>s$. The weights match exactly: the $Z_k$-dissipation in
$L_{k-1}$ has scale power $h^{2k-2}$, and that in $L_1$ has power $h^k$.

Here is the parameter argument. If $J_*>0$, set $z_k=J_k/J_*$ and
$\ell=L/J_*$. Using $z_k\le1$ in \eqref{eq:scalar-transfer} and descending
from $k=s$ gives a finite sum of bounds
\begin{equation}\label{eq:scalar-descent}
 z_k\le C_\G\sum_{r=k}^s\sum_{\theta\in\Theta_{kr}}
       (\ell/a_r)^\theta,\qquad
 2^{-s}\le\theta\le\tfrac12.
\end{equation}
The finite sets $\Theta_{kr}$ depend only on the step. Equation
\eqref{eq:scalar-upper}, summed in $i$, gives
\[
 \ell\le C_\G\left[1+\sum_{j<k}a_j(z_k+\sqrt{z_k})\right].
\]
Take $a_1=1$ and, for $k\ge2$,
\begin{equation}\label{eq:hierarchy}
 a_k=M^{B^{k-2}}.
\end{equation}
Choose the integer $B$ sufficiently large compared with $2^s$ and the degrees
of all the prescribed polynomials. Inserting \eqref{eq:scalar-descent} gives
\[
 \ell\le C_\G\left[1+M^{-c_s}\sum_{\theta\in\Theta}\ell^\theta\right],
 \qquad 0<\theta<1,\quad c_s>0.
\]
For sufficiently large $M$, this gives $\ell\le C_\G$, independently of $M$.
Then \eqref{eq:scalar-descent} gives $z_k\le C_\G a_k^{-c_s}$, after reducing
$c_s$ if needed. Since $a_k=a_{k-1}^B$, increasing $B$ first and then $M$
makes all $P_k(a_{<k})z_k$ and $P_k(a_{<k})/a_k$ as small as prescribed.
Also $\sum_{k\ge2}z_k\le1/2$, so $J_*\le2J_1$ and
$L\le C_\G J_1$. The case $J_*=0$ is immediate. For $s=1$, there are no
higher layers and $L_1\le J_1$ directly.
\end{proof}

This lemma is an existence statement for a separated hierarchy, not a claim
that raw gradient energies have uniform bounds for every arbitrary unordered
choice of diffusion weights. Its uniform constant is what permits subsequent
weights to be made larger without restarting the estimates.

\section{The exact tensor lift and removal of zero-order moments}
Use one index $\alpha$ for the full homogeneous basis and write $d_\alpha$ for
its degree. For a symmetric tensor $M$ on $\mathfrak g$, set
\[
 \cH M=\sum_{\alpha,\beta}Z_\alpha Z_\beta M_{\alpha\beta},\qquad
 C_\alpha=\ad_{Z_\alpha},\qquad
 \cC_\alpha M=C_\alpha M+MC_\alpha^T.
\]
A direct expansion gives
\begin{equation}\label{eq:first-intertwine}
 [Z_\alpha,\cH]=\cH\cC_\alpha.
\end{equation}
Therefore the exact tensor scale generator is
\begin{align}
 \cD_h^{\rm ten}
 &=\frac12\sum_\alpha a_{d_\alpha}h^{d_\alpha-1}
                         (Z_\alpha+\cC_\alpha)^2\notag\\
 &=\cA_h+\sum_\alpha a_{d_\alpha}h^{d_\alpha-1}
                  \cC_\alpha Z_\alpha+\cV_h,\label{eq:tensorgen}\\
 \cV_h&=\frac12\sum_\alpha a_{d_\alpha}h^{d_\alpha-1}\cC_\alpha^2.
\end{align}
Let $A_H$ be the horizontal embedding of $A$, and solve
\[
 N_h=\cD_h^{\rm ten}N,\qquad N(0)=mA_H.
\]
Equation \eqref{eq:first-intertwine} implies
\begin{equation}\label{eq:physical-lift}
 \cA_h\cH=\cH\cD_h^{\rm ten},\qquad
 S_h\cH(mA_H)=\cH N,\qquad \rho_t=\cH N.
\end{equation}
There is no residual commutator in the physical equation.

Define a finite-dimensional moment evolution
\begin{equation}\label{eq:moment}
 \Mom_h'=\cV_h\Mom_h,\qquad \Mom_0=I,
 \qquad \widehat N=\Mom_h^{-1}N.
\end{equation}
Each $C_\alpha$ raises degree, so $\Mom_h$ and $\Mom_h^{-1}$ are finite
polynomials in $h$ and the weights. This removes the zero-order forcing:
\begin{equation}\label{eq:renormalized-N}
 \widehat N_h=\cA_h\widehat N+
 \sum_\alpha a_{d_\alpha}h^{d_\alpha-1}
     \Mom_h^{-1}\cC_\alpha\Mom_h Z_\alpha\widehat N.
\end{equation}

The moment map preserves the positive-semidefinite cone. Indeed,
\[
 e^{\tau\cC_\alpha^2/2}M
 =\mathbb E\big[e^{\sqrt\tau\,\xi C_\alpha}M
                     e^{\sqrt\tau\,\xi C_\alpha^T}\big],
 \qquad \xi\sim N(0,1),
\]
and products and time limits preserve this property. Put
\[
 D_h^{\rm gr}=\operatorname{diag}(h^{(d_\alpha-1)/2}).
\]
For horizontal input $M_H$,
\begin{equation}\label{eq:moment-scaling}
 \Mom_h(M_H)=D_h^{\rm gr}\Mom_1(M_H)D_h^{\rm gr}.
\end{equation}
The $V_k\times V_k$ block of $\Mom_1(I_H)$ is bounded by a polynomial in
$a_1,\ldots,a_{k-1}$. Its horizontal block is exactly $I_r$.

\section{The layerwise tensor energy with strict weight independence}
For an unordered tensor block $\alpha=(i,k)$, $i\le k$, put
\[
 n_\alpha=i+k-2,\qquad
 B_\alpha=h^{-n_\alpha/2}\widehat N_{ik}/\rho.
\]
The horizontal block is
\[
 B_0=B_{11}=S_h(Am)/\rho,\qquad I\le B_0\le EI.
\]
It has the exact common-flow estimate
\begin{equation}\label{eq:base-F}
 F_0(h)=\int\rho^q|B_0|^2,\qquad F_0'=-\cD_q(B_0),\qquad
 \int_0^1\cD_q(B_0)\dd h\le rE^2Q(t,0).
\end{equation}

Order the remaining blocks by increasing total degree, refining ties in any
way. In \eqref{eq:renormalized-N}, a source into $B_\alpha$ has the form
\begin{equation}\label{eq:block-source}
 c_{\alpha\beta j}(a_{<k})\,h^{j/2-1}
       \big(Z_jB_\beta+B_\beta Z_jg\big),
 \qquad n_\beta<n_\alpha,\quad\max\beta\le k,\quad j<k.
\end{equation}
The coefficients are fixed linear maps and polynomials in weights strictly
below $k$. To verify the power of $h$, a conjugated term with degree increment
$n_\alpha-n_\beta$ has, before normalization, coefficient
$h^{(n_\alpha-n_\beta+j)/2-1}$. Normalization leaves $h^{j/2-1}$.
The restriction $j<k$ follows because $\ad_{Z_j}$ raises a nonzero input
index by $j$, so its output index is at least $j+1$. The same argument applies
to every factor in the moment maps.

In addition to the scalar quotient generator, $B_\alpha$ has the drift
$-n_\alpha B_\alpha/(2h)$. Set $F_\alpha=\int\rho^q|B_\alpha|^2$.
Its exact energy identity is
\begin{equation}\label{eq:block-energy}
 F_\alpha'+\cD_q(B_\alpha)+\frac{n_\alpha}{h}F_\alpha
 =2\sum_{\beta,j}h^{j/2-1}\int\rho^q
  \ip{B_\alpha}{c_{\alpha\beta j}(Z_jB_\beta+B_\beta Z_jg)}.
\end{equation}
The normalized nonhorizontal blocks have zero boundary value at $h=0$ for
smooth data: the moment forcing has been removed, and each remaining source
has at least one spatial derivative and an additional positive power
$h^{j/2}$ after time integration. Thus $F_\alpha(0)=0$.

\begin{lemma}\label{lem:tensor-energy}
For every $\alpha=(i,k)\ne(1,1)$ there is a polynomial
$K_\alpha(a_1,\ldots,a_{k-1})$ with nonnegative coefficients such that
\begin{equation}\label{eq:tensor-independent}
 F_\alpha(1)+\int_0^1\left[\cD_q(B_\alpha)+h^{-1}F_\alpha\right]\dd h
 \le K_\alpha(a_{<k})E^2\left[Q(t,0)+\sum_{j<k}J_j\right].
\end{equation}
The constants are independent of $a_k,\ldots,a_s$, of $q\in(1,3/2]$, and
of all coefficient derivatives.
\end{lemma}

\begin{proof}
For a nonhorizontal input block $\beta$, integrate its source in
\eqref{eq:block-energy} by parts:
\begin{equation}\label{eq:key-ibp}
 \int\rho^q\ip{B_\alpha}{c(Z_jB_\beta+B_\beta Z_jg)}
 =-\int\rho^q\ip{Z_jB_\alpha+\eps B_\alpha Z_jg}{cB_\beta}.
\end{equation}
Using \eqref{eq:second-square} and Young's inequality, this is bounded, after
including $2h^{j/2-1}$, by a small part of $\cD_q(B_\alpha)$ plus
\[
 C\norm c^2 a_j^{-1}h^{-1}F_\beta.
\]
For the horizontal input block, instead use $|B_0|\le\sqrt r E$ and the
positive damping $n_\alpha F_\alpha/h$. Its source is bounded by a small
part of that damping plus
\[
 C\norm c^2h^{j-1}\int\rho^q
          (|Z_jB_0|^2+rE^2|Z_jg|^2).
\]
Upon integration, \eqref{eq:base-F} controls the first term and the second is
$C\norm c^2rE^2J_j$. All $j$ occurring here are below $k$. Absorb the chosen
small parts in \eqref{eq:block-energy}, integrate, and induct in the strictly
increasing total degree. Since $a_j\ge1$, inverse $a_j$ factors may be
bounded by one. Squaring the finitely many source coefficients and multiplying
by previously obtained constants produces a polynomial only in $a_{<k}$.
This proves \eqref{eq:tensor-independent}.
\end{proof}

Define the physical remainder
\begin{equation}\label{eq:remainder}
 \cE_h=N-\Mom_h((\rho B_0)_H),\qquad
 E_{ik}=h^{-(i+k-2)/2}(\cE_h)_{ik}/\rho.
\end{equation}
By the polynomial scaling of $\Mom_h$, each $E_{ik}$ is a finite linear
combination of lower normalized blocks. Thus, enlarging the polynomials,
\begin{equation}\label{eq:remainder-energy}
 \int_0^1\frac1h\int\rho^q|E_{ik}|^2\dd h
 \le K_k(a_{<k})E^2\left[Q(t,0)+\sum_{j<k}J_j\right].
\end{equation}
Again there is no dependence on current or higher weights in $K_k$.

\section{Engel: the correction in explicit coordinates}
For $[X,Y]=Z$, $[X,Z]=W$, $[Y,Z]=0$, with $W$ central, use
\[
 \cA_h=\tfrac12(X^2+Y^2)+\tfrac12a_2hZ^2+\tfrac12a_3h^2W^2.
\]
Write the full tensor as
\[
 N=\begin{pmatrix}
 H_{11}&H_{12}&u_1&v_1\\
 H_{12}&H_{22}&u_2&v_2\\
 u_1&u_2&z&r\\
 v_1&v_2&r&w
 \end{pmatrix}.
\]
Put $V_1=(\cC_X^2+\cC_Y^2)/2$ and $V_2=\cC_Z^2/2$. All products of total
raising degree at least six vanish. Hence
\[
 \Mom_h=I+hV_1+\tfrac12h^2(V_1^2+a_2V_2),\qquad
 \Mom_h^{-1}=I-hV_1+\tfrac12h^2(V_1^2-a_2V_2).
\]
The renormalized blocks are
\begin{align}
 \widehat H&=H,& \widehat u&=u,\notag\\
 \widehat v&=v-\tfrac h2(H_{12},H_{22})^T,&
 \widehat z&=z-h\tr H,\notag\\
 \widehat r&=r-\tfrac32h u_2,&
 \widehat w&=w-h(v_2+z)+\tfrac14h^2(2H_{11}+3H_{22})
                         -\tfrac12a_2h^2H_{11}.\label{eq:engel-corrections}
\end{align}
Normalize $\widehat u,\widehat v,\widehat z,\widehat r,\widehat w$ by,
respectively, $\sqrt h\rho,h\rho,h\rho,h^{3/2}\rho,h^2\rho$.
The layer-two energy has a constant independent of both $a_2,a_3$; the
layer-three energy has a constant depending on $a_2$ but independent of $a_3$.
This is exactly the strict dependency required for choosing the weights.

For horizontal constant input $A$, the positive deterministic moment tensor is
\[
 \Mom_h(A_H)=
 \begin{pmatrix}
 A_{11}&A_{12}&0&hA_{12}/2\\
 A_{12}&A_{22}&0&hA_{22}/2\\
 0&0&h\tr A&0\\
 hA_{12}/2&hA_{22}/2&0&h^2(2a_2A_{11}+2A_{11}+3A_{22})/4
 \end{pmatrix}.
\]
Thus all renormalized nonhorizontal blocks vanish for spatially constant
$m,A$. The nonintegrable geometric variance in an uncorrected running-scale
energy has been removed, not estimated in absolute value.

\section{Physical-time coupling and closure}
Put $\cU(t)=\int_0^1Q(t,h)\dd h$. Split the exact physical flux as
\[
 \rho_t=\cH\big(\Mom_h((\rho B_0)_H)\big)+\cH\cE_h.
\]
We now estimate both parts. All parameter choices in this section precede the
choice of $\eps$.

\subsection{The positive moment part}
In the scaled covector frame $\widehat g_j=h^{(j-1)/2}\nabla_jg$, put
$M=\Mom_1((B_0)_H)$. Positivity of $\Mom_1$ gives $M\ge0$,
$M_{11}=B_0\ge I$, and
\[
 M_{kk}\le E P_k(a_{<k})I
\]
for fixed polynomials $P_k$. Positivity and a block Cauchy--Schwarz estimate give
\begin{equation}\label{eq:moment-good}
 \ip{M\widehat g}{\widehat g}
 \ge\tfrac12|\nabla_Hg|^2
       -C_\G E\sum_{k\ge2}P_k(a_{<k})h^{k-1}|\nabla_kg|^2.
\end{equation}
After one integration by parts, the other term differentiates $B_0$ only.
The norm of the fixed linear map $\Mom_1$, \eqref{eq:base-F}, and
$J_*\le2J_1$ give
\begin{equation}\label{eq:base-derivative-error}
 \left|\int_0^1\!\int\rho^q
 \sum_{\alpha,\beta}\widehat g_\alpha
 \big[\Mom_1((\widehat Z_\beta B_0)_H)\big]_{\alpha\beta}\right|
 \le\tfrac1{16}J_1+C_aQ(t,0),
\end{equation}
where $\widehat Z_\beta=h^{(d_\beta-1)/2}Z_\beta$ and $C_a$ is finite and
independent of $q$ and coefficient derivatives. Dependence of $C_a$ on all
weights is harmless: it multiplies $q\eps Q(t,0)$, and $\eps$ is chosen last.
Lemma~\ref{lem:scalar} makes the higher-layer sum in
\eqref{eq:moment-good} at most $J_1/16$.

\subsection{The corrected tensor remainder}
This is where two integrations by parts, rather than a pointwise tensor bound,
are decisive. For a block with degrees $i\le k$, its contribution to $Q_t$ is a
finite sum of terms
\begin{equation}\label{eq:two-ibp}
 q\eps\int\rho^q h^{(i+k-2)/2}E_{ik}
          \big(Z_k(Z_ig)+\eps(Z_kg)(Z_ig)\big),
\end{equation}
plus terms with $[Z_i,Z_k]g$. Indeed, interchange the two derivatives in the
symmetrized Hessian so that the higher-layer derivative is outside; the
resulting bracket has degree $i+k$.

Let $\mathfrak E_{ik}=\int_0^1h^{-1}\int\rho^q|E_{ik}|^2\dd h$.
Equation \eqref{eq:second-square} and Cauchy--Schwarz in space and scale give
\begin{equation}\label{eq:flux-estimate}
 \left|\int_0^1\text{the $(i,k)$ remainder in $Q_t$}\dd h\right|
 \le C_\G q\eps\sqrt{\mathfrak E_{ik}}
       \left(\sqrt{L_i/a_k}+\sqrt{J_{i+k}}\right),
\end{equation}
where $J_{i+k}=0$ if $i+k>s$. The weights are exact:
$h^{(i+k-2)/2}=h^{-1/2}h^{(i+k-1)/2}$.

Using \eqref{eq:remainder-energy} and the scalar lemma, the first part of
\eqref{eq:flux-estimate} is bounded by
\[
 C_\G q\eps E\sqrt{K_k(a_{<k})/a_k}
             \sqrt{Q(t,0)+J_1}\sqrt{J_1}.
\]
Make $E^2K_k(a_{<k})/a_k$ small. For the bracket terms prescribe in
Lemma~\ref{lem:scalar} the finitely many polynomials
\[
 \widetilde P_l(a_{<l})=E^2\sum_{i+k=l}K_k(a_{<k}).
\]
Their dependency is strictly lower because $k<i+k=l$. Choose their weighted
sum times $J_l$ small. The lemma and its simultaneous ratio conditions give
one finite choice of weights satisfying all these requirements. Consequently
all remainder terms together are bounded by
\begin{equation}\label{eq:small-flux}
 q\eps\left(\tfrac1{16}J_1+C_aQ(t,0)\right),
\end{equation}
after enlarging $C_a$.

\subsection{The closed entropy inequality}
Combining \eqref{eq:moment-good}--\eqref{eq:small-flux} gives
\begin{equation}\label{eq:time-J1}
 \cU'(t)\le-\tfrac14q\eps J_1+C_aq\eps Q(t,0).
\end{equation}
All constants and smoothing weights have already been fixed. From
\eqref{eq:Qgap} and $J_*\le2J_1$,
\[
 Q(t,0)-Q(t,1)\le q\eps C_{\rm gap}J_1,
 \qquad C_{\rm gap}=\max_ja_j.
\]
Choose
\begin{equation}\label{eq:eps-choice}
 0<\eps\le\min\left\{\tfrac12,\frac1D,
               \frac1{12C_{\rm gap}(1+C_a)}\right\}.
\end{equation}
Then
\begin{equation}\label{eq:closed}
 \boxed{\displaystyle
 \cU'(t)+\frac1{8C_{\rm gap}}Q(t,0)
       \le\frac1{4C_{\rm gap}}Q(t,1).}
\end{equation}
The coefficient error is absorbed only now, by choosing $q-1$ small. No
smoothing weight was chosen using a coefficient derivative or $1/(q-1)$.

\section{Adjoint estimate and the elliptic maximum estimate}\label{sec:analytic}
By mass preservation and \eqref{eq:kernel},
\[
 Q(t,h)\le C_\G^\eps h^{-D\eps/2}.
\]
Integrate \eqref{eq:closed} on $0<t<1$, use $\cU(1)\ge0$, and note that
$D\eps/2<1$. It follows that
\begin{equation}\label{eq:adjoint}
 \norm{P_t^*m_0}_{L^q((0,1)\times\G)}
 \le C(\G,E,q)\norm{m_0}_1.
\end{equation}
This is first proved for smooth $A$ and good positive initial data, uniformly
in all coefficient derivatives, and then for nonnegative $L^1$ initial data by
approximation and adjoint $L^1$ contraction. The same proof holds for every
smaller $\eps>0$ than the fixed choice in \eqref{eq:eps-choice}.

\subsection{Analytic justification for the smooth calculation}
For smooth $A$ equal to $I$ outside a compact set, the operator $L_A$ is a
smooth horizontal diffusion satisfying H\"ormander's condition. Its conservative
positive semigroup and smooth solutions are standard smooth-coefficient
inputs; no quantitative constant from such results is used in
\eqref{eq:closed}. For example, smooth coefficients are included in the
fundamental-solution theory in \cite{BLU,Faini}.

Conservativity can also be checked directly: in exponential coordinates the
coefficients of the equation for a coordinate in layer $k$ are polynomials in
strictly lower-layer coordinates, multiplied by bounded entries of $A$ or its
square root. The first-layer process has bounded covariance. Induction in $k$
and stopped moment estimates give finite moments on every finite time interval
and rule out explosion. This argument does not bound a derivative of $A$.

For the a priori calculation one may take a smooth positive initial density
with sufficiently fast polynomial decay and with its derivatives bounded by
an integrable polynomial weight. Smooth-coefficient kernel bounds and local
regularity, or the Duhamel formula against the fixed sub-Laplacian outside the
compact perturbation, justify spatial cutoffs and all integrations by parts.
The scalar and finite triangular tensor scale equations have the same
admissible decay. At $h=0$, the renormalized higher blocks tend to zero at a
positive power of $h$ for smooth inputs. Coefficient derivatives are used only
qualitatively for these endpoint limits. None enters the bounds proved before
passing to the limits. Initial densities in this class approximate arbitrary
nonnegative $L^1$ data. These whole-group arguments avoid any lattice issue.

\subsection{Resolvent duality and comparison}
For smooth nonnegative compactly supported $a$, duality in \eqref{eq:adjoint}
gives
\[
 \left\|\int_0^1P_ta\dd t\right\|_\infty\le C\norm a_{p_0},
 \qquad p_0=1+1/\eps.
\]
Positivity and $L^\infty$ contraction give, by splitting the time integral into
unit intervals,
\begin{equation}\label{eq:resolvent}
 v=(1-L_A)^{-1}a=\int_0^\infty e^{-t}P_ta\dd t\ge0,
 \qquad\norm v_\infty\le C_R\norm a_{p_0},\qquad L_Av=v-a.
\end{equation}
The solution is smooth by smooth hypoellipticity. Let $x_H$ denote first-layer
coordinates. On $B_1$ choose
\[
 b_0=(K-|x_H|^2)/(2r)>0.
\]
Then $L_Ab_0=-\tr A/r\le-1$. The usual interior-maximum argument applied to
$w-\tau b_0$ proves comparison on every open subset of $B_1$, without boundary
regularity: at an interior maximum the full Euclidean gradient vanishes and
$L_A(w-\tau b_0)\le0$.

\subsection{Measurable coefficients and the active set}
After subtracting the upper boundary value, fix $\delta>0$ and
$O_\delta=\{u>\delta\}\Subset\Omega$. Choose $\chi\in C_c^\infty(\Omega)$,
$0\le\chi\le1$, equal to one near $\overline O_\delta$ and supported in
$\{u>\delta/2\}$. Extend and mollify $A$ to obtain smooth $A_\nu$ equal to
$I$ outside a fixed compact set, with $I\le A_\nu\le EI$ and local convergence
in every finite $L^p$. Set $a_\nu=\chi(-L_{A_\nu}u)^+$. Then
\begin{equation}\label{eq:rough-limit}
 \norm{a_\nu}_{p_0}
 \le\norm{f^-\mathbf1_{\{u>\delta/2\}}}_{p_0}+o(1).
\end{equation}
The error on the fixed compact support is bounded by
$|A_\nu-A|\,|D_H^{2,*}u|$; this tends to zero in $L^{p_0}$.

Approximate each continuous compactly supported $a_\nu$ uniformly from above
by smooth nonnegative functions supported in a fixed slightly larger compact
set. Let $v=(1-L_{A_\nu})^{-1}a$ and $M=\norm v_\infty$. On $O_\delta$,
\[
 L_{A_\nu}(u-\delta-v-Mb_0)\ge-a-(v-a)+M\ge0,
\]
and the boundary values are nonpositive. Comparison and
\eqref{eq:resolvent} give
\[
 \sup_{O_\delta}(u-\delta)\le(1+\norm{b_0}_\infty)C_R\norm a_{p_0}.
\]
Pass through the smooth majorant approximation, then $\nu\to\infty$, then
$\delta\downarrow0$. This proves the normalized maximum estimate with the
active-set indicator. It uses no existence or uniqueness assertion for a
rough-coefficient resolvent or adjoint evolution.

Finally use $\Phi(g)=g_0\delta_Rg$, $\bar A=\lambda^{-1}A\circ\Phi$,
$\bar u=u\circ\Phi$. The forcing rescales as
\[
 \norm{\lambda^{-1}R^2f\circ\Phi}_p
 =\lambda^{-1}R^{2-D/p}\norm f_p.
\]
For $p\ge p_0$, use H\"older's inequality and $|\Omega|\le C_\G R^D$.
This proves Theorem~\ref{thm:abp}.

\section{Summary of the noncircular dependencies}
At tensor layer $k$, \eqref{eq:tensor-independent} and
\eqref{eq:remainder-energy} depend only on $a_1,\ldots,a_{k-1}$. The current
weight $a_k$ supplies the denominator in \eqref{eq:flux-estimate}; higher
weights remain available for the bracket term of degree $i+k$. The scalar
lemma handles any finite list of these strictly lower-dependent polynomials.
The coefficients of $Q(t,0)$ may depend on all weights, because $q-1$ is
chosen only after the weights have been fixed. This separates geometric
absorption from coefficient absorption and closes the scale entropy.

The argument uses neither an assumed rough-coefficient Gaussian kernel nor a
pointwise full-rank bound on the transported tensor. The replacement for the
kinetic Gaussian is the scalar all-layer kernel \eqref{eq:kernel}, its exact
tensor lift \eqref{eq:tensorgen}, and the finite moment correction
\eqref{eq:moment}.

\begin{thebibliography}{9}
\bibitem{Kinetic} \emph{Kinetic ABP through adjoint estimate}, supplied
unpublished manuscript, \texttt{kinetic ABP-adjoint.tex}.
\bibitem{Heisenberg} \emph{A short analytic proof of a Heisenberg ABP estimate},
supplied unpublished manuscript, \texttt{heisenberg-abp.tex}.
\bibitem{GTHeat} N. Garofalo and G. Tralli, \emph{A universal heat semigroup
characterisation of Sobolev and BV spaces in Carnot groups},
arXiv:2205.04574, 2022.
\bibitem{BLU} A. Bonfiglioli, E. Lanconelli and F. Uguzzoni,
\emph{Fundamental solutions for non-divergence form operators on stratified
groups}, Trans. Amer. Math. Soc. \textbf{356} (2004), 2709--2737.
\bibitem{Faini} M. Faini, \emph{Non-divergence evolution operators modeled on
H\"ormander vector fields with Dini continuous coefficients}, Mathematics in
Engineering \textbf{8} (2026), 477--517, doi:10.3934/mine.2026015.
\end{thebibliography}
\end{document}
